\documentclass{beamer}
\usetheme{Madrid}
\usecolortheme{dolphin}
\definecolor{charcoal}{RGB}{54,69,79}
\setbeamercolor{structure}{fg=charcoal}
\setbeamercolor{palette primary}{bg=charcoal,fg=white}
\setbeamercolor{palette secondary}{bg=charcoal!90,fg=white}
\setbeamercolor{palette tertiary}{bg=charcoal!80,fg=white}
\setbeamercolor{palette quaternary}{bg=charcoal!70,fg=white}
\setbeamercolor{frametitle}{bg=charcoal,fg=white}
\setbeamercolor{title}{bg=charcoal,fg=white}
\usepackage{hyperref}
\usepackage{graphicx}
\usepackage{booktabs}
\usepackage{listings}
\usepackage{color}
\definecolor{dkgreen}{rgb}{0,0.6,0}
\definecolor{gray}{rgb}{0.5,0.5,0.5}
\definecolor{mauve}{rgb}{0.58,0,0.82}
\lstset{
  language=Java,
  aboveskip=3mm,
  belowskip=3mm,
  showstringspaces=false,
  columns=flexible,
  basicstyle={\small\ttfamily},
  numbers=none,
  keywordstyle=\color{blue},
  commentstyle=\color{dkgreen},
  stringstyle=\color{mauve},
  breaklines=true,
  breakatwhitespace=true,
  tabsize=2
}
\title[Game Project Kickoff]{Game Project Kickoff}
\subtitle{COP2250: Java Programming}
\author{Kevin Pyatt, Ph.D.}
\institute{State College of Florida \\ Pyatt Labs}
\date{Week 12}
\begin{document}
\begin{frame}
  \titlepage
\end{frame}
\begin{frame}{What You're Building}
\textbf{Terminal Game Project --- team of 2--3 students}
\vspace{0.3cm}
\begin{itemize}
  \item Original interactive game written entirely in Java
  \item Runs in the terminal --- no GUI required
  \item You choose the game type
  \item Delivered in stages over the next 4 weeks
\end{itemize}
\vspace{0.3cm}
\textit{You choose the game. You design the architecture. You write every line.}
\end{frame}
\begin{frame}{Teams and HedgeDoc}
\textbf{Your team workspace is already set up.}
\vspace{0.3cm}
\begin{tabular}{ll}
\toprule
\textbf{Team} & \textbf{Workspace} \\
\midrule
RPG & Linked on Project Overview page \\
Maria \& Vel & Linked on Project Overview page \\
Jayden \& Draden & Linked on Project Overview page \\
Kayde, Diego, Ryan \& Edgar & Linked on Project Overview page \\
\bottomrule
\end{tabular}
\vspace{0.3cm}
Open now: \texttt{pyattlabs.github.io/scf-java-labs/\#/java-game/README}
\end{frame}
\begin{frame}{Project Timeline}
\begin{tabular}{ll}
\toprule
\textbf{Week} & \textbf{Milestone} \\
\midrule
Week 12 & Stages 1, 2, 3 --- Design, Architecture, Charter \\
Week 13 & Sprint 1 --- Core classes running \\
Week 14 & Sprint 2 --- Inheritance and object relationships \\
Week 15 & Sprint 3 --- Feature complete + Playtest \\
Week 16 & Final Demo + Individual Defense \\
\bottomrule
\end{tabular}
\vspace{0.3cm}
\textit{Design documents are 30\% of your project grade. Do not skip them.}
\end{frame}
\begin{frame}{Stage 1 --- Game Design Document}
\textbf{No code. Just the concept.}
\vspace{0.2cm}
\begin{enumerate}
  \item Game title and one-sentence pitch
  \item Game type and genre
  \item Core mechanics --- what does the player do?
  \item Win condition, lose condition, or end state
  \item Main entities in your game
  \item What actions can the player take?
\end{enumerate}
\vspace{0.3cm}
\textbf{Deliverable:} Posted to your HedgeDoc and submitted via Canvas.
\end{frame}
\begin{frame}{Stage 2 --- Architecture Document}
\textbf{Translate the GDD into class design.}
\vspace{0.2cm}
\begin{itemize}
  \item Class hierarchy diagram (hand-drawn, ASCII, or tool)
  \item Description of each class: responsibility, fields, methods
  \item Which classes are abstract and why
  \item Which interfaces you plan to implement and why
\end{itemize}
\vspace{0.3cm}
\textit{What are the things in your game? Each thing is probably a class.}
\end{frame}
\begin{frame}{Stage 3 --- Team Charter}
\textbf{Team agreements in writing. Required before Sprint 1.}
\vspace{0.2cm}
\begin{itemize}
  \item Who owns which classes
  \item How you communicate outside of class
  \item How you handle disagreements
  \item Signed acknowledgment from all members
\end{itemize}
\vspace{0.3cm}
\textit{No charter = no sprint credit.}
\end{frame}
\begin{frame}{OOP Requirements}
\begin{tabular}{ll}
\toprule
\textbf{Requirement} & \textbf{Minimum} \\
\midrule
Distinct classes & 4 or more \\
Abstract class & 1, with 2+ concrete subclasses \\
Method override & At least 1 \\
Composition & At least 1 \\
Interface & At least 1 \\
Exception handling & Required \\
Game loop & Required \\
End condition & Required \\
GUI reflection & Required from every member \\
\bottomrule
\end{tabular}
\end{frame}
\begin{frame}[fragile]{Sample Class Hierarchy}
\begin{lstlisting}
GameEngine
|-- World
|   |-- Room
|-- Character (abstract)
|   |-- Player
|   |-- NPC
|-- Item (abstract)
    |-- Weapon
    |-- Key
\end{lstlisting}
\vspace{0.3cm}
\begin{itemize}
  \item \texttt{Character} is abstract --- Player and NPC share behavior
  \item \texttt{Item} is abstract --- all items have a name and effect
  \item \texttt{GameEngine} composes \texttt{World} and \texttt{Player}
\end{itemize}
\end{frame}
\begin{frame}{Game Ideas If You're Stuck}
\begin{tabular}{ll}
\toprule
\textbf{Game} & \textbf{Core Classes} \\
\midrule
Blackjack & Deck, Card, Hand, Player, Dealer \\
Quiz Engine & QuizBank, Question, Player, ScoreTracker \\
Dungeon Crawler & Room, Enemy, Weapon, Player, GameEngine \\
Inventory Sim & Item, Warehouse, Order, Customer \\
\bottomrule
\end{tabular}
\vspace{0.3cm}
\textit{Simple to start. Depth is always rewarded. Get it running first.}
\end{frame}
\begin{frame}{Today's Work}
\textbf{By end of class:}
\vspace{0.2cm}
\begin{enumerate}
  \item Open your team HedgeDoc
  \item Complete Stage 1: Game Design Document
  \item Sketch Stage 2: class hierarchy draft
  \item Begin Stage 3: Team Charter
\end{enumerate}
\vspace{0.3cm}
\textbf{Demo Requirement:} Every team member must be able to explain\\
any part of the design document when called on.\\
\vspace{0.3cm}
\textit{You cannot fake understanding in front of a live terminal.}
\end{frame}
\end{document}
